% Source: MA-GY7033 Fall 2026 Homework 5.pdf, p. 2, Problem 8(a--e).
% Preserve the supplied superscript vector coordinates and column convention.
\begin{exercise}
  \label{ex:hw05:permutation-matrices}
  \solutionlink{hw05:permutation-matrices}
  Given a permutation \(\sigma\in S_n\), the \emph{permutation matrix}
  \(P_\sigma\) is the matrix whose \(k\)-th column is
  \(c_k=e_{\sigma(k)}\), i.e.,
  \[
    P_\sigma=\bigl[e_{\sigma(1)}\mid\cdots\mid e_{\sigma(n)}\bigr].
  \]
  For example, \(n=3\) and \(\sigma\) is the permutation given by
  \[
    \begin{aligned}
      \sigma(1)&=2\\
      \sigma(2)&=1\\
      \sigma(3)&=3,
    \end{aligned}
  \]
  then
  \[
    \begin{aligned}
      P_\sigma
        &=\bigl[e_{\sigma(1)}\mid e_{\sigma(2)}\mid e_{\sigma(3)}\bigr]\\
        &=\bigl[e_2\mid e_1\mid e_3\bigr]\\
        &=\begin{bmatrix}0&1&0\\1&0&0\\0&0&1\end{bmatrix}.
    \end{aligned}
  \]
  \begin{enumerate}[(a)][2]
    \item Show that a matrix is a permutation matrix if and only if every
      entry is \(0\) or \(1\) and there is exactly one \(1\) in each row
      and each column.
    \item Show that if
      \[
        M=\bigl[c_1\mid\cdots\mid c_n\bigr],
      \]
      then
      \[
        MP_\sigma=\bigl[c_{\sigma(1)}\mid\cdots\mid c_{\sigma(n)}\bigr],
      \]
    \item Given a permutation \(\sigma\) and
      \(v=(v^1,\ldots,v^n)\in\R^n\), describe the linear map defined
      by \(P_\sigma\). In other words, what is \(P_\sigma v\)?
    \item Show that, for any two permutations \(\sigma\) and \(\tau\),
      \[
        P_{\sigma\circ\tau}=P_\sigma P_\tau
      \]
      by showing that for any \(k\in\{1,\ldots,n\}\),
      \[
        P_{\sigma\circ\tau}e_k=P_\sigma P_\tau e_k.
      \]
    \item Show that \(\det(P_\sigma)=\epsilon(\sigma)\).
  \end{enumerate}
\end{exercise}
